% se puede agregar la opción [english] para 
%  memorias o tesis en inglés (borrando el archivo .aux)
\documentclass[hyphens]{umemoria} 

\depto{Departamento de Física}
\author{Diego Antonio Román Cortés}
\title{Redes multiorbitales basadas en moléculas fotónicas}

% incluir ambos comandos para una doble titulación
%  o quitar el comando que no aplica
%\memoria{Ingenier[o/a?] Civil en ???}
\tesis{Magíster en Ciencias, mención en Física}
%\tesis{Doctor en ???} % incluir solo este comando para doctorados

% puede haber varios profesores guía seperados por coma;
% pero si es una memoria, solo puede haber un profesor guía
\guia{Rodrigo Andrés Vicencio Poblete} 

% puede haber varios profesores co-guía seperados por coma;
% pero si es una memoria, el profesor co-guía será el primer
% integrante de la comisión
%\coguia{Nombre Completo Co-Guía} % incluir en caso de co-guía de *tesis*

%\cotutela{Nombre Institución} % incluir en caso de cotutela

\comision{Jaime Anguita García, Pedro Orellana}

\auspicio{los proyectos Instituto Milenio para la Investigación en Óptica (MIRO) ICN17 012, Fondecyt Regular 1191205 y
1231313} % incluir en caso de recibir financiamiento

% tiene que ser el año en que se da el examen de título/grado (defensa)
%\anho{2021} % incluir solo para reemplazar el año actual

\usepackage{lipsum}
\usepackage{txfonts}
\usepackage{graphicx}
%\usepackage[spanish, es-nodecimaldot]{babel}
\usepackage{float}
\usepackage[figurename=Figura]{caption}
\usepackage[numbers, sort&compress]{natbib}
%\usepackage[hyphens]{url}
\usepackage{indentfirst}
\usepackage{hyperref}
\usepackage{xcolor}
\usepackage{ragged2e}

\begin{document}

\frontmatter
\maketitle

\begin{resumen}
\lipsum[1-4]
\end{resumen}

% opcional: incluir para tesis en inglés;
%  en este caso hay que tener el resumen y abstract
%   en ambos idiomas
%\begin{abstract}
%\lipsum[1-4]
%\end{abstract}

\begin{dedicatoria}
Una dedicatoria corta.
\end{dedicatoria}

\begin{thanks}
\lipsum[1-2]
\end{thanks}

\tableofcontents
\listoftables % opcional
\listoffigures % opcional

\mainmatter

\chapter{Introducción}

Durante la última década, varios Premios Nobel en Física han estado estrechamente ligados a la óptica \cite{nobel}: por la generación de pulsos ultracortos de luz —femtosegundos \cite{femto1} y luego attosegundos \cite{atto1, atto2, atto3}—, por experimentos pioneros con fotones entrelazados \cite{photons1, photons2, photons3}, por el desarrollo de pinzas ópticas \cite{opticaltweezers} y por la invención de LEDs de alta eficiencia \cite{led1, led2, led3}. Estos avances fundamentales han impulsado aplicaciones industriales, biomédicas, en telecomunicaciones e incluso en defensa. Una aplicación cotidiana destacada es la fibra óptica, que actúa como guía de onda para la luz y constituye hoy el principal medio de transmisión de datos a nivel global \cite{fibra2, fibra}.

Muchos de estos desarrollos se han visto potenciados por la técnica de \textit{escritura láser de femtosegundos}, que permite fabricar redes fotónicas tridimensionales de guías de onda con precisión micrométrica. Esta tecnología ha habilitado la realización experimental de redes con geometrías complejas y grados de libertad artificiales, y ha sido utilizada para estudiar fenómenos como oscilaciones de Bloch \cite{BlochOsci}, localización de Anderson \cite{Anderson}, dinámica de bandas planas \cite{lieb1, lieb2, artificialFB, FBdynamics}, fases topológicas \cite{obstopo, obsfloquet, topo1dphoto, toporusos}, formación de solitones no lineales \cite{discretesolitons} e incluso la propagación de luz cuántica \cite{qed, squeezed, topoquantum}.

Dentro de este marco, el presente trabajo se enfoca en el estudio de redes fotónicas \textit{multiorbitales}, en las que los modos transversales de la luz —en particular los modos \( P \), de simetría antisimétrica— actúan como grados de libertad internos que pueden aprovecharse para enriquecer la dinámica óptica. A diferencia de los sistemas convencionales, donde se propaga solamente el modo fundamental (\( S \)), las redes multiorbitales permiten emular estructuras más complejas, como acoplamientos interorbitales o interferencias moduladas por la geometría.

Uno de los fenómenos experimentales que motiva este estudio es el \textit{ángulo de invisibilidad} \cite{Pmodecoupling}, descrito en el Capítulo~\ref{cap:invisibility}. En ciertas condiciones geométricas, los modos \( P \) se desacoplan completamente unos de otros, permaneciendo ópticamente \textit{invisibles}. Este fenómeno pone de manifiesto que la simetría transversal de los modos no solo es relevante, sino que puede ser usada activamente como herramienta de control. Este principio guía la exploración de nuevas formas de localización, filtrado y direccionamiento de luz en dispositivos fotónicos integrados.

Luego, se utiliza el concepto de \textit{acoplamiento interorbital SP}, un mecanismo que permite hibridar modos \( S \) y \( P \) ajustando finamente sus constantes de propagación a través del control de la potencia del láser de escritura \cite{interorbital}. Este acoplamiento da lugar a una dinámica rica, especialmente cuando se implementa en redes periódicas.

En el Capítulo~\ref{cap:molecules}, este principio se aplica a la construcción de moléculas fotónicas —pares de guías fuertemente acopladas— cuyas combinaciones simétrica y antisimétrica funcionan como orbitales efectivos. Sobre esta base, se implementa experimentalmente el modelo \textit{SP-SSH} \cite{SPSSH}, una extensión multiorbital del modelo de Su-Schrieffer-Heeger \citep{ssh}. Este sistema presenta una \textit{doble transición de fase topológica}, determinada por dos parámetros independientes: la dimerización geométrica y el desintonizado de las constantes de propagación entre los $S$ y $P$. La existencia de tres fases distintas —con dos, cuatro o ningún modo de borde— se analiza mediante la polarización de bulto y técnicas espectrales.

Ambas líneas —la invisibilidad angular y el acoplamiento SP— se articulan bajo una visión común: aprovechar la \textit{estructura geométrica de los modos guiados} como un nuevo recurso físico en fotónica.

\vspace{1em}

Este trabajo se organiza de la siguiente manera:

El Capítulo~\ref{cap:teo} introduce el formalismo teórico utilizado a lo largo de la tesis, incluyendo los principios de propagación modal, acoplamiento en redes discretas y conceptos topológicos asociados a bandas ópticas.

En el Capítulo~\ref{cap:num} se presentan las herramientas numéricas utilizadas. Se describe el método de Expansión en Modos Normales (EME), que permite estudiar el comportamiento estacionario de la luz, junto con otras metodologías como el método de propagación de haces (BPM) y la teoría de modos acoplados (CMT).

El Capítulo~\ref{cap:exp} detalla la técnica experimental de escritura láser, con énfasis en los parámetros relevantes para diseñar redes multiorbitales tridimensionales.

En el Capítulo~\ref{cap:invisibility} se estudia el fenómeno de invisibilidad modal de los modos \( P \), caracterizando experimentalmente la geometría crítica que conduce al desacoplamiento óptico de estos modos. Esta propiedad se pone a prueba en una red tipo panal de abeja. Para lograr una concordancia cuantitativa entre la teoría y el experimento, se introduce un modelo extendido que incorpora tanto la \textit{no ortogonalidad} entre los modos como \textit{acoplamientos de largo alcance} entre guías no adyacentes.

El Capítulo~\ref{cap:molecules} analiza la formación de moléculas fotónicas y la ingeniería de sus modos internos. A partir de estas, se implementa y caracteriza una red SP-SSH, incluyendo el estudio de sus fases topológicas mediante simulaciones y análisis espectrales.

Finalmente, el Capítulo~\ref{cap:conclu} presenta las conclusiones del trabajo y discute posibles extensiones hacia otras plataformas fotónicas.




\chapter{Marco teórico}

\section{Modos normales del campo electromagnético en un medio dieléctrico}

Esta tesis estudia el comportamiento de luz láser de baja potencia (1 mW de potencia de salida) propagada en guías de onda dieléctricas escritas dentro de una muestra de borosilicato. Es por ello que se supone un medio lineal no magnético libre de fuentes de carga y de corriente. Las ecuaciones de Maxwell en éste régimen son:
\begin{align}
	\nabla\cdot\textbf{D} &= 0, \label{eqn:gauss}
	\\	
	\nabla\times\textbf{E} &= -\frac{\partial \textbf{B}}{\partial t}, \label{eqn:faraday-lenz}
	\\	
	\nabla\cdot\textbf{B} &= 0,
	\\	
	\nabla\times\textbf{B} &= \mu_0\frac{\partial \textbf{D}}{\partial t}, \label{eqn:ampere-maxwell}
\end{align}
donde \textbf{E}, \textbf{B}, $\textbf{D}=\varepsilon\textbf{E}$ son los campos eléctrico, campo magnético y campo desplazamiento eléctrico, respectivamente. Las guías de onda son invariantes en la dirección de propagación $z$, por lo que el índice de refracción $n=\sqrt{\varepsilon/\varepsilon_0}$ dependerá de las coordenadas transversales al eje óptico, es decir, $n \equiv n(x,y) = n_0 + \Delta n(x,y)$, con $n_0=1.47$ el índice de refracción del borosilicato y $\Delta n \sim 10^{-5}-10^{-3}$ el contraste de las guías de onda.

Aplicando rotor por la izquierda a la ecuación (\ref{eqn:faraday-lenz}), usando la ecuación (\ref{eqn:ampere-maxwell}) y asumiendo una solución temporal harmónica proporcional a $e^{i\omega t}$ se tiene:

\begin{align}
	\nabla\times\nabla\times\textbf{E} &= -\frac{\partial}{\partial t}(\nabla\times\textbf{B}) = -\frac{\partial}{\partial t}\left(\mu_0\frac{\partial \textbf{D}}{\partial t}\right) = -\frac{n^2}{c^2}\frac{\partial^2 \textbf{E}}{\partial t^2} = n^2k_0^2 \textbf{E}, \label{eqn:rotordoble}
\end{align}
donde $k_0 \equiv \omega/c$ es el número de onda en el vacío. Notemos que, por identidad de cálculo vectorial, se tiene que $\nabla\times\nabla\times\textbf{E} = \nabla(\nabla\cdot\textbf{E}) - \nabla^2\textbf{E}$, y usando la ley de Gauss (\ref{eqn:gauss}) se deduce que $\nabla\cdot \textbf{E} = -\nabla(n^2)\cdot\textbf{E}/n^2$.

Con esto, 

\begin{equation}
	(\nabla^2  + k_0^2n^2)\textbf{E} = -\nabla\left(\textbf{E} \cdot \frac{\nabla n^2}{n^2}\right) \approx 0, \label{eqn:helmholz}
\end{equation}

donde el término de la derecha se desprecia usando la aproximación de guiaje débil, pues el contraste $\Delta n$ es pequeño. Si la estructura de guías de onda no varía en la dirección $z$, el campo eléctrico se puede expandir en sus modos normales como $\textbf{E}(\textbf{r}) = \sum_k \textbf{E}_k(x, y) e^{i\beta_k z}$, por lo que el laplaciano de la ecuación (\ref{eqn:helmholz}) se puede separar como $\nabla^2 \equiv \nabla_\perp^2 + \frac{\partial^2}{\partial z^2}$, por lo que la ecuación (\ref{eqn:helmholz}) se expande como:

\begin{align}
	(\nabla^2  + k_0^2n^2) \textbf{E}(\textbf{r}) &= \left(\nabla_\perp^2 + \frac{\partial^2}{\partial z^2} + k_0^2n^2\right)\sum_k \textbf{E}_k(x, y)  e^{i\beta_k z} \nonumber
\\	
	&= \sum_k \left[ e^{i\beta_k z} \nabla_\perp^2 \textbf{E}_k + -\beta_k^2\textbf{E}_k e^{i\beta_k z} + k_0^2n^2 \textbf{E}_k  e^{i\beta_k z}\right]
\nonumber	
	\\	
	&= \sum_k \left[  \nabla_\perp^2  + (k_0^2n^2-\beta_k^2) \right]\textbf{E}_k  e^{i\beta_k z}
	\nonumber	
	\\
	&=
	0
	\nonumber
	\\
	\implies
	& \left[  \nabla_\perp^2  + k_0^2n^2(x,y) \right]\textbf{E}_k(x,y)  = \beta_k^2 \textbf{E}_k(x,y), \quad \forall k. \label{eqn:eigenfield}
\end{align}

Notemos que la ecuación (\ref{eqn:eigenfield}) es un problema de autovalores $\beta_k^2$ y autofunciones $\textbf{E}_k(x,y)$. En principio, la forma espacial del índice de refracción $n(x, y)$ puede ser arbitraria siempre que se satisfaga la condición de guiaje débil de la ecuación (\ref{eqn:helmholz}). 

\section{Soluciones analíticas para un contraste $n(x)$ rectangular}

\section{Teoría de modos acoplados}
	Para el estudio de redes fotónicas, es conveniente utilizar herramientas similares a las de la Física del Sólido en lo que respecta a potenciales periódicos. En particular, se puede suponer que los modos guiados de una guía de onda están fuertemente ligados a ella (enlace fuerte o \textit{Tight Binding}), incluso en presencia de otras guías de onda. Es decir, se supondrá que el $k$-ésimo modo de la $m$-ésima guía de onda satisface para toda distancia de propagación $z$ la ecuación (\ref{eqn:eigenfield}), donde el índice de refracción total se puede descomponer en una suma periódica de guías de onda $n^2(\textbf{r}) = \sum_{m} n^2_m(\textbf{r})$. Entonces, despomponiendo el campo eléctrico total de la forma $\textbf{E}(\textbf{r}) = \sum_{k, m} \textbf{E}_{k, m}(x, y) a_{k, m}(z) e^{i\beta_{k, m} z}$ y reemplazando en la ecuación (\ref{eqn:helmholz}) se tiene:

\begin{align}
	(\nabla^2  + k_0^2n^2) \textbf{E}(\textbf{r}) &= \left(\nabla_\perp^2 + \frac{\partial^2}{\partial z^2} + k_0^2n^2 \right)\sum_{k, m} \textbf{E}_{k, m}(x, y) a_{k, m}(z) e^{i\beta_{k, m} z}
	\nonumber
	\\
	&= \sum_{k, m} \left[a_{k, m} e^{i\beta_{k, m} z} \left(\nabla_\perp^2 +k_0^2n^2 \right)\textbf{E}_{k, m} + \textbf{E}_{k, m}\frac{d^2}{d z^2}a_{k, m} e^{i\beta_{k, m} z}\right]
	\nonumber	
	\\
	&= \sum_{k, m} \left[a_{k, m}  \left(\nabla_\perp^2 +k_0^2n^2 -\beta_{k,m}^2 \right) + \frac{d^2 a_{k, m}}{d z^2}  +2i\beta_{k,m}\frac{d a_{k, m}}{d z} \right]e^{i\beta_{k, m} z}\textbf{E}_{k, m}
		\nonumber	
	\\
	&= \sum_{k, m} \left[a_{k, m}  k_0^2\sum_{m'\neq m} n^2_m(\textbf{r})  +2i\beta_{k,m}\frac{d a_{k, m}}{d z} \right]e^{i\beta_{k, m} z}\textbf{E}_{k, m}
	\nonumber	
	\\
	&= 0
	\nonumber	
	\\
	\implies &
	  \int_{-\infty}^{+\infty}\int_{-\infty}^{+\infty} \sum_{k, m} \left[a_{k, m}  k_0^2\sum_{m'\neq m} n^2_m(\textbf{r})  +2i\beta_{k,m}\frac{d a_{k, m}}{d z} \right]e^{i\beta_{k, m} z}\textbf{E}_{k, m} \cdot \textbf{E}_{k', m'}^* dxdy = 0
	\label{eqn:CMT1}
\end{align}

	


\section{Métodos numéricos}

A partir de la ecuación (\ref{eqn:rotordoble}) y separando el campo eléctrico $\textbf{E}(\textbf{r}) = \textbf{E}_1(\textbf{r}) e^{-i \phi_1 (\textbf{r})}$ en una envolvente lenta $\textbf{E}_1(\textbf{r})$ y una fase rápidamente oscilante $\phi_1 (\textbf{r})$ ,

\begin{align}
\nabla\times[\nabla\times(\textbf{E}_1 e^{-i \phi_1 })] &=  n^2k_0^2 \textbf{E}_1 e^{-i \phi_1}
\\
\nabla\times[ e^{-i \phi_1} \times (\nabla \times \textbf{E}_1) + \nabla(e^{-i \phi_1})\times \textbf{E}_1] &=  n^2k_0^2 \textbf{E}_1 e^{-i \phi_1}
\\
\nabla\times[ e^{-i \phi_1}  (\nabla  - i  \textbf{k} _1)\times\textbf{E}_1  ] &=  n^2k_0^2 \textbf{E}_1 e^{-i \phi_1}
\\
e^{-i \phi_1}  \nabla\times[ (\nabla  - i  \textbf{k} _1)\times\textbf{E}_1] + \nabla(e^{-i \phi_1})\times  (\nabla  - i  \textbf{k} _1)\times\textbf{E}_1  &=  n^2k_0^2 \textbf{E}_1 e^{-i \phi_1}
\\
e^{-i \phi_1}  \nabla\times[ (\nabla  - i  \textbf{k} _1)\times\textbf{E}_1] - i\textbf{k}_1 e^{-i \phi_1}\times  (\nabla  - i  \textbf{k} _1)\times\textbf{E}_1  &=  n^2k_0^2 \textbf{E}_1 e^{-i \phi_1}
\\
	(\nabla-i\textbf{k}_1)\times((\nabla-i\textbf{k}_1)\times \textbf{E}_1) &= n^2k_0^2 \textbf{E}_1,
	 \label{eqn:comsol}
\end{align}
con $\textbf{k}_1 = \nabla\phi_1(\textbf{r})$.
La ecuación (\ref{eqn:comsol}) se puede resolver vía el software comercial COMSOL \textit{Multiphysics} mediante elementos finitos más gruesos que los que se tendría que usar a partir de la ecuación (\ref{eqn:rotordoble}) debido a la separación entre envolvente y fase, sin haber hecho aproximación alguna. Por otro lado, utilizando la aproximación paraxial escogiendo el eje $z$  como dirección de propagación y seleccionando la polarización horizontal del campo, es posible simplificar la ecuación (\ref{eqn:comsol}) y llegar a la formulación de un método numérico escalar conocido como \textit{Beam Propagation Method}, utilizado ampliamente en esta área de investigación \cite{bics, interorbital, OAMCaging, vortex, bpm}. Éste consiste en resolver numéricamente la ecuación
\begin{equation}
	2in_0k_0\frac{\partial}{\partial z}\psi(x,y,z) = \nabla_\perp^2 \psi (x,y,z) + \left(n^2-n_0^2\right)k_0^2 \psi (x,y,z), \label{eqn:paraxial}
\end{equation}
con $\psi(x,y,z)$, $n_0$ , $n(x,y)$ y $k_0$ la envolvente de la componente horizontal campo eléctrico, el índice de refracción del material, el índice de refracción inducido y el número de onda en el vacío, respectivamente. Aplicando teoría acoplada de modos \cite{coupledmodetheory} a la ecuación (\ref{eqn:paraxial}), con el objetivo de describir de forma discreta una red fotónica, es posible derivar las llamadas ecuaciones discretas tipo Schrödinger \cite{discretesolitons, artificialFB, FBdynamics}
\begin{equation}
	-i\frac{\partial u_{\vec{n}} }{\partial z} = \beta_{\vec{n}}u_{\vec{n}} + \sum_{\vec{m}\neq\vec{n}} C_{\vec{n},\vec{m}}u_{\vec{m}}, \label{eqn:CMT}
\end{equation}
con $u_{\vec{n}}, \beta_{\vec{n}}$ y $C_{\vec{n}, \vec{m}}$ la envolvente normalizada del campo eléctrico, la constante de propagación normalizada y las constantes de acoplamiento entre los modos de las guías en las posiciones de la red $\vec{n}$ y $\vec{m}$, respectivamente.
\chapter{Tercero}
\lipsum[50-60]
\chapter{Conclusiones}

\lipsum[130-132]
\begin{figure}
	\centering
	\includegraphics[scale=.2]{imagenes/fcfm.pdf}
	\caption{Logo de la Facultad}
	\label{logofcfm}
\end{figure}

\lipsum[133-134]

\begin{table}
	\centering
	\caption{Tabla 1}
	\label{tabla:1}
	\begin{tabular}{|c|c|r|}
		\hline
		\textbf{Campo 1} & \textbf{Campo 2} & \textbf{Num} \\\hline
		Valor 1a & Valor 2a & 3\\
		Valor 1b & Valor 2b & 3\\
		\hline
	\end{tabular}

\end{table}

\lipsum[135]


% ver https://www.overleaf.com/learn/latex/Glossaries
% \input{glosario.tex} % opcional

%\nocite{*}
\bibliographystyle{unsrtnat}
\bibliography{citations}

% opcional ...
\begin{appendices}
\section{Ortogonalidad de los modos normales \label{sec:orto}}

La ortogonalidad de los modos normales $\textbf{E}_\nu^\perp$ puede demostrarse utilizando el convenio de Einstein en la ecuación (\ref{eqn:eigenfield}) para simplificar la notación:
\begin{equation}
	T_{ij} E^\nu_j = \beta_\nu^2 E^\nu_i, \label{eqn:eigentensorial}
\end{equation}
donde el operador $T_{ij}$ actúa como $T_{ij}E_j \equiv \delta_{ij}\left(\nabla_\perp^2 + k_0^2n^2\right)E_j + \partial_i \left(E_j \partial_j\ln(n^2)\right).$

Para demostrar la ortogonalidad de los modos $E_i$, se puede demostrar equivalentemente que $T_{ij}$ es hermítico:
\begin{align}
	\iint \left(T_{ij} E_j^\mu\right)^* E_i^\nu \,dA = \iint \left(E_i^\mu\right)^* \left(T_{ij} E_j^\nu\right) \,dA.
\end{align}
Esta propiedad se sigue de la identidad:
\begin{align*}
	&\left(T_{ji} E_i^\mu\right)^* E_j^\nu - \left(E_i^\mu\right)^* \left(T_{ij} E_j^\nu\right) =  \nabla_\perp \cdot \left[E_i^\nu \nabla_\perp \left(E_i^\mu\right)^* - \left(E_i^\mu\right)^* \nabla_\perp E_i^\nu\right]  + \left[E^\nu_j \partial_j \left(E_i^\mu\right)^* - \left(E_j^\mu\right)^* \partial_j E_i^\nu\right] \partial_i \ln(n^2).
\end{align*}
Al integrar sobre el plano transversal el término de divergencia se anula para modos guiados (que decaen en el infinito), pero el segundo término desaparece sólo si $\nabla n^2 = \textbf{0}$. Es decir, el operador $T_{ij}$ sólo es Hermítico en la condición de guiaje débil.

El sistema completo es hermítico, como se demuestra partiendo de la ecuación (\ref{eqn:rotordoble}):
\begin{align*}
	\iiint_V \left(\nabla\times\nabla\times\textbf{E}^\nu\right) \cdot \left(\textbf{E}^\mu\right)^* dV &= \iiint_V n^2\frac{\omega_\nu^2}{c^2} \textbf{E}^\nu \cdot \left(\textbf{E}^\mu\right)^* dV \\
	\iiint_V \left[\nabla\times\nabla\times\left(\textbf{E}^\mu\right)^*\right] \cdot \textbf{E}^\nu dV &= \iiint_V n^2 \frac{\omega_\mu^2}{c^2} \left(\textbf{E}^\mu\right)^* \cdot \textbf{E}^\nu dV
\end{align*}

Restando ambas ecuaciones y aplicando el teorema de divergencia:
\begin{align*}
	\frac{\omega_\nu^2 - \omega_\mu^2}{c^2} \iiint_V n^2 \textbf{E}^\nu \cdot \left(\textbf{E}^\mu\right)^* dV &= \oiint_{\partial V} \left[ \cdots \right] \cdot \hat{\textbf{n}}\,dA \\
	&\xrightarrow{\partial V \to \infty} 0
\end{align*}
Se obtienen así dos posibilidades: o los modos son degenerados ($\omega_\nu^2 = \omega_\mu^2$), o se tiene que $\iiint n^2 \textbf{E}^\nu \cdot \left(\textbf{E}^\mu\right)^* dV = 0$.
\end{appendices}
\end{document}